% ============================================================================
% 01-introduction.tex — DRAFTED IN FULL
% ============================================================================
\section{Introduction}\label{sec:intro}

Ap\'ery's proof of the irrationality of $\zeta(3)$ placed a handful of integer
sequences --- defined by three-term recurrences with polynomial coefficients,
yet integral against all apparent odds --- at the center of a circle of ideas
connecting holonomic differential equations, modular forms, and the geometry
of families of algebraic surfaces. The six \emph{sporadic} Ap\'ery-like
sequences of order 2, catalogued by Zagier, and their order-3 companions have
since been objects of sustained experimental and theoretical attention.
Cooper~\cite{Cooper2012} identified three further sporadic sequences, commonly
denoted $s_7$, $s_{10}$ and $s_{18}$, satisfying order-3 recurrences;
Gorodetsky~\cite{Gorodetsky2023} gave closed forms and placed all three inside
a single four-parameter recurrence template. Almkvist and van
Straten~\cite{AlmkvistvanStraten2021} exhibited explicit projective K3 models
whose periods realize the corresponding order-3 operators.

This paper studies the differential operators attached to Cooper's template
from three angles --- structural, arithmetic, and lattice-theoretic --- with a
methodological commitment that we regard as part of the contribution:
\emph{every claim is labeled by how it is known}. Specifically, each statement
in this paper carries exactly one of the following statuses, and the label is
part of the statement's presentation, not a footnote:

\begin{enumerate}
\item[\textbf{(K)}] \emph{Kernel-checked}: proved in Lean~4 and verified by
  its proof kernel against the pinned \texttt{mathlib}
  library~\cite{deMouraUllrich2021,mathlibCPP2020,mathlib4}, with the axiom
  base disclosed in \S\ref{sec:formalization};
\item[\textbf{(E)}] \emph{Exact}: established by exact symbolic computation
  over $\Q$ (no floating point), reproducible by a named checker script in the
  accompanying repositories, each shipped with negative controls
  (\S\ref{sec:reproducibility});
\item[\textbf{(N)}] \emph{Numerically supported}: derived from
  high-precision numerics followed by exact recognition and exact
  post-verification; stated as a \emph{Numerical Claim}, never as a theorem;
\item[\textbf{(C)}] \emph{Cited}: a classical result of the literature, used
  as stated, with the precise statement and location given.
\end{enumerate}

\subsection{Results}

Write $\thetaop = z\,\frac{d}{dz}$. Cooper's template, in the normalization
of~\cite{Gorodetsky2023}, is the family of order-3 operators
\begin{equation}\label{eq:cooper-template}
  \Lthree(a,b,c,d)
  \;=\;
  \thetaop^3 - z\,(2\thetaop+1)\bigl(a\thetaop^2 + a\thetaop + b\bigr)
  + z^2\,\bigl(c(\thetaop+1)^3 + d(\thetaop+1)\bigr),
  \qquad (a,b,c,d)\in\Q^4,
\end{equation}
with $(a,b,c,d) = (13,4,-27,3)$, $(6,2,-64,4)$, $(14,6,192,-12)$ for $s_7$,
$s_{10}$, $s_{18}$ respectively \cite{Cooper2012,Gorodetsky2023}.

Our first result is a uniform structure theorem
(Theorem~\ref{thm:sym2-template}, status \textbf{(K)}): for \emph{every}
parameter choice, $\Lthree(a,b,c,d)$ is, in a division-free
$\thetaop$-form sense made precise in \S\ref{sec:sym2}, the product
$P_2\cdot\Symtwo(\Ltwo)$ of its own leading $\thetaop$-coefficient with the
symmetric square of an order-2 operator $\Ltwo$ whose coefficients are
\emph{determined} by the template --- the partner is solved out of the
identities, leaving a single nontrivial constraint, which holds identically.
Independently, the Almkvist--van Straten symmetric-square criterion $W\equiv 0$
holds identically on the template in cleared-denominator form
(Theorem~\ref{thm:avs}, status \textbf{(K)}). Both facts are single
kernel-checked proofs covering $s_7$, $s_{10}$, $s_{18}$ --- and any future
instance of the template --- simultaneously by specialization. To our
knowledge this is the first kernel-checked verification of a symmetric-square
structure theorem for Picard--Fuchs-type operators in any proof assistant.

The partner operator makes an arithmetic distinction visible
(\S\ref{sec:integrality}): among Cooper's three sporadic sequences, only the
$s_7$ partner has an integral holomorphic solution --- the sequence
$1, 2, 22, 336, 6006, \dots$ (OEIS A279619~\cite{OEIS}). Non-integrality for
$s_{10}$ and $s_{18}$ is kernel-checked by finite witness (status
\textbf{(K)}); integrality for $s_7$ is a theorem of
O'Brien~\cite{OBrien2016} (status \textbf{(C)}), whose applicability to our
partner is kernel-checked at the level of recurrences. We isolate the
mechanism: the modular coordinate $X_7 = \eta_1^3\eta_7^3/z_7^3$ is a
\emph{normalized} integral uniformizer, and normalization --- not mere
integrality of $\eta$-quotients --- is the load-bearing property (status
\textbf{(E)}).

Third (\S\ref{sec:irreducibility}, status \textbf{(E)}): the $s_7$ and
$s_{10}$ partners $\Ltwo$ are irreducible, by a denominator obstruction ---
every attainable residue sum of a hyperexponential logarithmic derivative lies
in $\tfrac12\Z$, while the exponents of $\Ltwo$ at infinity are
$\{\tfrac13,\tfrac23\}$ (resp.\ $\{\tfrac38,\tfrac58\}$) --- and their
monodromy is not dihedral, because the doubled indicial root at $0$ forces a
logarithmic solution and no nontrivial unipotent element lies in the
normalizer of a maximal torus of $\SL_2$. Consequently
$\Lthree = \Symtwo(\Ltwo)$ is irreducible and is the \emph{minimal-order}
annihilator of its holomorphic solution. All steps are exact over $\Q$.

Finally (\S\ref{sec:lattice}), we compute the monodromy representation of the
$s_7$ partner by numerical analytic continuation at 60-digit working
precision and recognize the entries exactly; the largest recognition residual
is ${\approx}10^{-59}$ against an acceptance gate of $10^{-35}$, and the
recognized matrices pass exact structural post-verification (involutivity,
the cusp relation at infinity, existence and uniqueness up to scale of a
joint invariant form, orbit-lattice closure). This step is numerics-backed
and is reported as Numerical Claim~\ref{nc:monodromy} (status \textbf{(N)}),
not as a theorem. Taking the recognized matrices as input, an exact
integer-arithmetic pipeline then derives the joint monodromy-invariant
lattice in its primitive even scaling and proves --- unconditionally, as a
statement about the exhibited Gram matrix (Proposition~\ref{prop:usplit},
status \textbf{(E)}) --- that it is isometric over $\Z$ to
$\U\oplus\lat{14}$: Gram matrix
$\left[\begin{smallmatrix}0&0&-1\\0&14&0\\-1&0&0\end{smallmatrix}\right]$,
determinant $-14$, signature $(2,1)$, discriminant form $\Z/14$ with
$q=\tfrac1{14}$. The isometry is realized by an explicit $\GL_3(\Z)$ base
change and has been verified by \emph{two independently written} exact
implementations; the identical pipeline applied to the $s_{10}$ family
derives determinant $-20$ and $\U\oplus\lat{20}$, so the computation
discriminates between levels. This matches on the nose the lattice
$(M_7)^{\perp} = \U\oplus\lat{14}$ of Dolgachev's $M_n$-polarized K3
framework~\cite{Dolgachev1996}, in which the Picard--Fuchs operator is a
symmetric square by a theorem of Doran~\cite{Doran2000}; \S\ref{sec:framework}
assembles the evidence that the $s_7$ family is $M_7$-polarized, and states
precisely what remains unproven --- in particular, the identification of the
computed monodromy lattice with the transcendental lattice of the family is
presented conditionally, not as a theorem.

\subsection{What is new}

None of the framework results are ours: symmetric-square structure for
$M_n$-polarized families is Doran's theorem, and $\smash{(M_n)^{\perp} =
\U\oplus\lat{2n}}$ is Dolgachev's computation. The contributions of this
paper are: (i) the uniform, kernel-checked template theorem with an
explicitly determined partner, where the literature treats candidates one at
a time and asserts existence rather than exhibiting the partner from the
parameters; (ii) the exact irreducibility/minimality proofs for these
specific operators; (iii) the sharpened integrality mechanism for $s_7$,
including the checked resolution of an ambiguity in the source's defining
equation; (iv) the explicit monodromy lattice of the $s_7$ family with an
exact, doubly-verified integral splitting $\U\oplus\lat{14}$ --- a
\emph{converse-direction} computation: rather than deducing invariants from a
polarization assumption, we compute the invariants of a concretely given
operator family and match them against the framework, in a situation where,
as Doran notes (\cite{Doran2000}, \S6), the classification of rank-19
lattice polarizations is not complete; and (v) the epistemic architecture
itself --- kernel-checked/exact/numerical/cited labeling with reproducible
artifacts and negative controls --- which we offer as a methodological
contribution to experimental mathematics.

\subsection{Organization}

\S\ref{sec:prelim} fixes notation. \S\ref{sec:operators} presents the two
operators for $s_7$ and their exact Riemann schemes. \S\ref{sec:sym2} states
and proves the template structure theorem. \S\ref{sec:irreducibility} proves
irreducibility and minimality. \S\ref{sec:lattice} presents the monodromy
computation and the integral splitting. \S\ref{sec:integrality} treats the
$s_7$ integrality mechanism. \S\ref{sec:framework} relates the results to the
Dolgachev--Doran framework and states the conditional statements and open
questions. \S\ref{sec:formalization} describes the Lean~4 formalization,
including its complete axiom inventory. \S\ref{sec:reproducibility} indexes
the artifacts needed to reproduce every number in this paper.

% ----------------------------------------------------------------------------
% Generated-by: Fable 5 (Stream 1 paper agent) | Verified-by: claims traced to
% paper/PLAN.md #3 rows 1-20 | Reviewed-by: pending T0 (Xavier)
% ----------------------------------------------------------------------------
